% ===========================================================================
% Chapitre 5 — Validation, sécurité et exploitation.
%
% Les résultats sont RÉELS et DATÉS (2026-09-22). Ils doivent être redatés
% après correction des trois échecs. Le vocabulaire d'état est contraint :
% « réussite supposée » n'en fait pas partie.
% ===========================================================================
\chapter{Validation, sécurité et exploitation}
\label{chap:validation}

\section{Introduction}
\label{sec:val_intro}

Ce chapitre présente la validation de PIQBIT. Il repose sur une distinction
que le reste du rapport applique systématiquement : un mécanisme
\emph{implémenté} n'est pas un mécanisme \emph{validé}. La première affirmation
se lit dans le code ; la seconde exige une exécution capturée et datée.

Cette distinction explique le vocabulaire d'états employé ici. « Réussite
capturée » signifie qu'une exécution datée existe. « Implémentation
confirmée » signifie que le mécanisme existe mais qu'aucune exécution n'a été
capturée. « Non vérifiée » signifie qu'aucune preuve n'existe. Ces trois états
ne sont pas interchangeables, et aucun n'est un synonyme atténué d'un autre.

\section{Stratégie de validation}
\label{sec:test_strategy}

Le tableau~\ref{tab:val_strategy} expose la stratégie retenue.

\begin{table}[htbp]
  \centering
  \caption{Stratégie de validation}
  \label{tab:val_strategy}
  \begin{tabularx}{\textwidth}{L{3.2cm} L{3.4cm} X}
    \toprule
    \textbf{Niveau} & \textbf{Moyen} & \textbf{Objet} \\
    \midrule
    Unitaire        & Suite de tests serveur & Règles métier, calcul des scores, transitions de statut \\
    Intégration     & Suite de tests sur base réelle & Persistance, autorisation, parcours multi-services \\
    Interface       & Build de production & Typage strict, budgets de taille, compilation \\
    Sécurité        & Tests fonctionnels dédiés & Authentification, verrouillage, contrôle d'accès \\
    Modèle          & Partitions indépendantes & Erreur de prédiction mesurée hors apprentissage \\
    Exécution       & Sondes et point de santé & Démarrage des conteneurs et disponibilité \\
    \bottomrule
  \end{tabularx}
\end{table}

\paragraph{Écart entre environnements, à déclarer.} La suite serveur s'exécute
sur une base embarquée en local, et sur le moteur cible — extension vectorielle
comprise — en intégration continue. Ces deux exécutions ne sont pas
équivalentes : un résultat obtenu sur la base embarquée ne vaut pas pour le
moteur cible, notamment sur les chemins vectoriels. Le rapport ne transpose
donc jamais un résultat de l'un à l'autre.

\section{Résultats de la campagne du 22 septembre 2026}
\label{sec:validation_results}

Le tableau~\ref{tab:validation_results} synthétise la campagne.

\begin{table}[htbp]
  \centering
  \caption{Résultats de validation capturés le 22 septembre 2026}
  \label{tab:validation_results}
  \begin{tabularx}{\textwidth}{L{3.4cm} L{4.1cm} X}
    \toprule
    \textbf{Validation} & \textbf{État de preuve} & \textbf{Résultat} \\
    \midrule
    Build de production de l'interface
      & \evstate{Réussite capturée}
      & Code de sortie 0 en 36,0 s ; bundle initial 1,25 Mo brut et
        246,82 ko transférés ; \textbf{quatre avertissements de budget}. \\
    Suite de tests serveur
      & \evstate{Échec capturé}
      & 594 réussis, \textbf{3 échoués}, 6 ignorés en 607,4 s ;
        code de sortie 1. \\
    Contrôle de dérive du schéma
      & \evstate{Implémentation confirmée}
      & Configuré en intégration continue ; aucune exécution capturée. \\
    Images et démarrage de la pile
      & \evstate{Implémentation confirmée}
      & Aucune construction ni démarrage capturés. \\
    Tests unitaires de l'interface
      & \evstate{Planifié}
      & Non exécutés dans cette campagne. \\
    Bout en bout, performance, résilience
      & \evstate{Planifié}
      & Aucun harnais, aucune mesure. \\
    Analyse statique
      & \evstate{Sans objet}
      & Aucun outil configuré dans le projet. \\
    \bottomrule
  \end{tabularx}
\end{table}

\pqscreenshot{ci_workflow_summary}
  {Exécution de la chaîne d'intégration continue : les deux travaux, leur
   durée et leur issue.}
  {fig:screen_ci}

\section{Analyse des échecs}
\label{sec:failure_analysis}

Trois tests ont échoué. Les trois sont analysés, car un échec non expliqué est
plus inquiétant qu'un échec compris.

\paragraph{E-01 — synchronisation d'agenda.} Attendu \texttt{200}, obtenu
\texttt{502}. Ce test s'exécute sur un poste où l'agenda externe est configuré,
et la substitution de l'appel réseau n'est pas empruntée par le chemin réel.
Il s'agit d'une \emph{défaillance du dispositif de test}, dépendante de
l'environnement : elle ne démontre ni un défaut du produit, ni son bon
fonctionnement. Le test doit être réécrit pour isoler l'appel externe.

\paragraph{E-02 — test antérieur au durcissement des autorisations.} Attendu
\texttt{200}, obtenu \texttt{403}. Le refus est \emph{conforme} à la politique
en vigueur ; c'est l'attente du test qui est périmée, écrite avant le
durcissement. Le test doit être réécrit avec un acteur autorisé. Ce n'est pas
un défaut du produit.

\paragraph{E-03 — incohérence d'autorisation.} Attendu \texttt{403}, obtenu
\texttt{200}. Celui-ci est différent des deux précédents : c'est un
\textbf{défaut réel}.

Un rôle que la politique entend exclure atteint une action qui engage
financièrement l'entreprise, parce que l'endpoint est gardé par une permission
\emph{plus large} que celle visée. L'écart n'est pas dans le moteur
d'autorisation — qui applique exactement ce qui lui est déclaré — mais entre
l'\emph{intention} de la politique et la \emph{déclaration} de la route.

Deux précisions s'imposent sur la portée de ce constat. D'une part, il provient
d'un test fonctionnel et non d'un audit de sécurité : on ne peut en tirer
aucune conclusion générale sur la sécurité du système. D'autre part, il doit
être corrigé puis suivi d'une \emph{nouvelle campagne datée} — le rapport final
devant refléter l'état livré et non un état intermédiaire.

\pqscreenshot{ci_test_failure_detail}
  {Détail des trois échecs de la campagne du 22 septembre 2026 dans le journal
   d'intégration continue.}
  {fig:screen_failures}

\section{Validation fonctionnelle par domaine}
\label{sec:functional_validation}

Le tableau~\ref{tab:functional_validation} récapitule la couverture par
domaine.

\begin{table}[htbp]
  \centering
  \caption{Couverture de validation fonctionnelle par domaine}
  \label{tab:functional_validation}
  \begin{tabularx}{\textwidth}{L{3.4cm} L{3.2cm} X}
    \toprule
    \textbf{Domaine} & \textbf{État} & \textbf{Observation} \\
    \midrule
    Authentification et 2FA & \evstate{Réussite capturée} & Verrouillage et égalisation de latence couverts par tests dédiés \\
    Autorisation            & \evstate{Échec capturé}     & Écart E-03 sur un endpoint sensible \\
    Offres et candidatures  & \evstate{Réussite capturée} & Contrôles de dépôt et cycle de vie couverts \\
    Analyse et notation     & \evstate{Réussite capturée} & Inclut le cas de non-régression de la surnotation \\
    Rapprochement sémantique& \evstate{Partielle}         & Seul le chemin de repli est exercé \\
    Assistant extractif     & \evstate{Réussite capturée} & Présence des sources vérifiée \\
    Entretiens              & \evstate{Échec capturé}     & Échec E-01, dépendant de l'environnement \\
    Test technique          & \evstate{Réussite capturée} & Chronométrage et correction couverts \\
    Salaire et négociation  & \evstate{Réussite capturée} & Métriques de modèle relues dans l'artefact daté \\
    Contrats et signature   & \evstate{Réussite capturée} & Gardes d'état et idempotence couvertes \\
    RH interne              & \evstate{Réussite capturée} & Couverture de base \\
    Rapports et notifications& \evstate{Implémentation confirmée} & Diffusion temps réel non capturée \\
    \bottomrule
  \end{tabularx}
\end{table}

\section{Qualité du modèle de prédiction salariale}
\label{sec:model_quality}

Le modèle livré a été sélectionné parmi huit candidats comparés sur cinq
partitions indépendantes. Le tableau~\ref{tab:model_metrics} rapporte les
métriques relues dans l'artefact daté du 10 septembre 2026.

\begin{table}[htbp]
  \centering
  \caption{Métriques du modèle de prédiction salariale retenu}
  \label{tab:model_metrics}
  \begin{tabularx}{\textwidth}{L{5.0cm} L{3.2cm} X}
    \toprule
    \textbf{Métrique} & \textbf{Valeur} & \textbf{Lecture} \\
    \midrule
    Erreur absolue moyenne (test) & 402 TND & Écart moyen entre salaire prédit et salaire réel, sur données non vues \\
    Coefficient de détermination  & 0,891   & Part de la variance expliquée par le modèle \\
    Modèles comparés              & 8       & Sélection sur l'erreur de test, non d'entraînement \\
    Partitions d'évaluation       & 5       & Partitions indépendantes, pour ne pas conclure sur un découpage favorable \\
    \bottomrule
  \end{tabularx}
\end{table}

Deux précautions de lecture. D'abord, une erreur absolue moyenne de 402 dinars
n'est pas une garantie individuelle : c'est une moyenne, et des profils
atypiques peuvent s'en écarter sensiblement. Ensuite, ces métriques sont
mesurées sur le jeu de données disponible ; leur transposition à une population
de candidats différente n'est pas démontrée. Le modèle produit donc une
\emph{indication} destinée à éclairer une négociation, non une valeur de
référence.

\section{Validation de sécurité}
\label{sec:security_validation}

Le tableau~\ref{tab:security_validation} distingue ce qui est couvert de ce
qui ne l'est pas.

\begin{table}[htbp]
  \centering
  \caption{Contrôles de sécurité et état de validation}
  \label{tab:security_validation}
  \begin{tabularx}{\textwidth}{L{3.8cm} L{3.2cm} X}
    \toprule
    \textbf{Contrôle} & \textbf{État} & \textbf{Moyen} \\
    \midrule
    Hachage des mots de passe      & \evstate{Réussite capturée} & Test fonctionnel dédié \\
    Verrouillage de compte         & \evstate{Réussite capturée} & Test fonctionnel dédié \\
    Égalisation du temps de réponse& \evstate{Réussite capturée} & Test fonctionnel dédié \\
    Second facteur                 & \evstate{Réussite capturée} & Test fonctionnel dédié \\
    Contrôle d'accès par rôle      & \evstate{Échec capturé}     & Écart E-03 \\
    Entropie des jetons publics    & \evstate{Implémentation confirmée} & Lecture du code \\
    Gardes d'état des liens publics& \evstate{Réussite capturée} & Tests de rejeu \\
    Analyse de dépendances         & \evstate{Sans objet}        & Aucun outil configuré \\
    Recherche de secrets           & \evstate{Sans objet}        & Aucun outil configuré \\
    Sécurité dynamique             & \evstate{Planifié}          & Aucun harnais \\
    Chiffrement du transport       & \evstate{Absent}            & Aucune terminaison TLS dans la composition fournie \\
    \bottomrule
  \end{tabularx}
\end{table}

\pqscreenshot{rbac_denied_403}
  {Refus d'autorisation lorsqu'un rôle non habilité vise une route réservée.}
  {fig:screen_rbac}

\pqscreenshot{2fa_totp_enrollment}
  {Activation du second facteur d'authentification. Le secret affiché est un
   secret jetable, produit pour la capture.}
  {fig:screen_2fa}

\section{Portes de qualité de la chaîne d'intégration}
\label{sec:ci_gates}

La figure~\ref{fig:delivery_flow} décrit la chaîne et, surtout, son point
d'arrêt.

\pqfigure{delivery_flow}
  {Flux de validation et point d'arrêt de la chaîne. Deux travaux parallèles
   alimentent une porte de qualité ; au-delà, tout est manuel et non capturé.}
  {fig:delivery_flow}

Le tableau~\ref{tab:ci_gates} distingue ce qui est \emph{configuré} de ce qui
est \emph{exécuté et daté}.

\begin{table}[htbp]
  \centering
  \caption{Portes de qualité : configuré contre exécuté}
  \label{tab:ci_gates}
  \begin{tabularx}{\textwidth}{L{4.2cm} L{2.8cm} X}
    \toprule
    \textbf{Porte} & \textbf{État} & \textbf{Observation} \\
    \midrule
    Tests serveur sur moteur cible  & Configuré et exécuté & Échec au 22 septembre 2026 \\
    Contrôle de dérive du schéma    & Configuré            & Exécution non capturée \\
    Build de production             & Configuré et exécuté & Réussi, avec avertissements de budget \\
    Tests unitaires de l'interface  & Configuré            & Non exécutés dans cette campagne \\
    Analyse statique                & Absent               & Aucun outil dans le projet \\
    Recherche de secrets            & Absent               & Aucun outil dans le projet \\
    \bottomrule
  \end{tabularx}
\end{table}

Au 22 septembre 2026, la porte de qualité était donc \textbf{fermée} : la suite
serveur échouait.

\section{Livraison, déploiement et exploitation}
\label{sec:delivery}

Le tableau~\ref{tab:ops_state} énonce l'état réel de l'exploitation.

\begin{table}[htbp]
  \centering
  \caption{Capacités d'exploitation : existant et manquant}
  \label{tab:ops_state}
  \begin{tabularx}{\textwidth}{L{4.0cm} L{2.4cm} X}
    \toprule
    \textbf{Capacité} & \textbf{État} & \textbf{Observation} \\
    \midrule
    Journalisation applicative & Présente & Configurée au démarrage, sans donnée personnelle de CV \\
    Point de contrôle de santé & Présent  & Interrogeable par les sondes \\
    Sondes de conteneur        & Présentes& Forcées sur la quatrième version du protocole \\
    Journal d'audit            & Présent  & Acteur, action, cible \\
    Métriques                  & Absentes & Aucune télémétrie implémentée \\
    Traces distribuées         & Absentes & — \\
    Tableau de bord d'exploitation & Absent & — \\
    Alertes                    & Absentes & — \\
    Publication d'artefacts    & Absente  & La chaîne ne publie rien \\
    Déploiement automatisé     & Absent   & Aucun environnement cible \\
    Procédure de retour arrière& Absente  & Non définie \\
    \bottomrule
  \end{tabularx}
\end{table}

Aucune observabilité de bout en bout ne peut donc être revendiquée. Les
capacités présentes permettent de constater qu'un service répond ; elles ne
permettent pas de mesurer son comportement ni d'être alerté de sa dégradation.
Combler cet écart suppose trois briques distinctes : un traçage distribué
normalisé~\cite{opentelemetry}, une collecte spécifique aux chemins reposant
sur des modèles de langage~\cite{langfuse}, et une évaluation systématique de
la qualité de récupération~\cite{es2024ragas}. Les deux dernières sont
nécessaires parce qu'une trace d'appel ordinaire indique la latence et le code
de statut, mais pas si la réponse était fondée.

\pqscreenshot{docker_compose_healthy}
  {Pile de conteneurs démarrée : les quatre services et leur état de santé.}
  {fig:screen_docker}

\pqscreenshot{health_endpoint}
  {Réponse du point de contrôle de santé de l'interface de programmation.}
  {fig:screen_health}

\section{Limites}
\label{sec:limitations}

Les limites de la solution livrée sont récapitulées dans le
tableau~\ref{tab:limitations}. Elles sont énoncées ici plutôt que dispersées,
afin qu'aucune ne puisse être manquée.

\begin{table}[htbp]
  \centering
  \caption{Limites constatées de la solution livrée}
  \label{tab:limitations}
  \begin{tabularx}{\textwidth}{L{1.3cm} L{3.2cm} X}
    \toprule
    \textbf{Réf.} & \textbf{Domaine} & \textbf{Limite} \\
    \midrule
    L-01 & Autorisation   & Écart E-03 : un rôle non habilité atteint une action sensible \\
    L-02 & Interface      & Budget de taille du bundle dépassé (NFR-08) \\
    L-03 & Sémantique     & Seul le chemin de repli est validé ; le chemin vectoriel natif ne l'est pas \\
    L-04 & Exploitation   & Aucune métrique, trace, alerte ni tableau de bord \\
    L-05 & Déploiement    & Aucun environnement cible, aucune terminaison TLS, aucun retour arrière \\
    L-06 & Performance    & Aucun test de charge ni de résilience \\
    L-07 & Accessibilité  & Aucune campagne conduite (NFR-11) \\
    L-08 & Disponibilité  & Non mesurée faute de production (NFR-10) \\
    L-09 & Sécurité       & Aucune analyse de dépendances, de secrets, ni sécurité dynamique \\
    L-10 & Modélisation   & Deux liens non contraints par la base : employés/comptes et négociation/candidat \\
    L-11 & Signature      & Preuve d'intention, non signature qualifiée au sens réglementaire \\
    L-12 & Modèle salarial& Métriques établies sur le jeu disponible ; transposition non démontrée \\
    \bottomrule
  \end{tabularx}
\end{table}

\section{Travaux futurs}
\label{sec:future_work}

\begin{enumerate}[leftmargin=*]
  \item Corriger l'écart E-03, réécrire les tests E-01 et E-02, puis produire
        une \emph{nouvelle campagne datée} qui reflète l'état livré.
  \item Réduire le bundle initial sous son budget, par découpage plus fin et
        réexamen des dépendances lourdes.
  \item Capturer une exécution du chemin vectoriel natif, afin de ne plus
        valider la seule voie de repli.
  \item Introduire une télémétrie minimale — métriques d'usage, latence des
        parcours critiques — sans laquelle aucune exploitation sérieuse n'est
        possible.
  \item Ajouter une terminaison TLS et définir un environnement cible avec
        procédure de retour arrière.
  \item Ajouter une analyse de dépendances et une recherche de secrets aux
        portes de qualité.
  \item Conduire une campagne d'accessibilité et un test de charge.
\end{enumerate}

\section{Conclusion du chapitre}
\label{sec:val_conclusion}

Ce chapitre a présenté la validation de PIQBIT avec des résultats datés du
22 septembre 2026. Le build de production réussit, avec un dépassement de
budget signalé ; la suite serveur échoue sur trois tests, dont un seul
constitue un défaut réel — l'écart d'autorisation E-03. Le modèle salarial
atteint une erreur absolue moyenne de 402 dinars sur données non vues, avec une
méthodologie de sélection justifiée.

Douze limites sont consignées. Les déclarer relève de la même discipline que
de rapporter les réussites : un rapport qui tairait le dépassement de budget ou
l'absence d'environnement cible serait contredit par sa propre annexe de
validation.

La conclusion générale dresse le bilan objectif par objectif et propose les
perspectives d'extension de la plateforme.
